% ============================================================================
% Solusi Bab 7: Metode Simpleks (bentuk kamus)
% Sumber buku: part1-linear-programming/ch06-simplex/simplex-basis-driven.tex
% Penomoran latihan mengikuti urutan \begin{ex} dalam file tersebut (17 latihan).
% Semua pivot dan optimum diverifikasi dengan aljabar basis pecahan eksak dan
% scipy.optimize.linprog, Juli 2026.
% ============================================================================
\section{Bab 7: Metode Simpleks}

Tujuh belas latihan dalam bab ini mengikuti alur bab: enam soal pemanasan tentang
bentuk standar dan kamus awal (7.1 sampai 7.6), tujuh penyelesaian simpleks inti
dalam bentuk kamus termasuk dua masalah Big-M (7.7 sampai 7.13), tiga
soal konseptual tentang kondisi optimalitas, geometri solusi basis
layak, dan uji rasio (7.14 sampai 7.16), serta satu soal tantangan tentang
degenerasi akibat hasil imbang pada uji rasio (7.17). Setiap pivot di bawah ini diverifikasi melalui
aljabar basis dengan pecahan eksak dan setiap optimum dikonfirmasi dengan pemecah
PL. Solusi yang dicetak dalam Solusi Terpilih buku (7.5, 7.6, 7.7,
7.8, 7.9, 7.14, 7.17) direproduksi di sini dengan sedikit penyuntingan agar
manual ini mandiri.

\exsol{7.1}{Bentuk standar dan kamus pertama}{1}

\textbf{(a)} Tambahkan variabel slack $s_1, s_2 \geq 0$:
\[
\begin{aligned}
\max \quad & 4x + 5y \\
\st \quad & x + 2y + s_1 = 12, \\
& 3x + 2y + s_2 = 24, \\
& x, y, s_1, s_2 \geq 0.
\end{aligned}
\]

\textbf{(b)} Menyelesaikan setiap persamaan terhadap variabel slack-nya memberikan kamus awal
pada basis slack $\{s_1, s_2\}$:
\[
\begin{aligned}
\max \quad & z = 0 + 4x + 5y \\
\st \quad & s_1 = 12 - x - 2y, \\
& s_2 = 24 - 3x - 2y, \\
& x, y, s_1, s_2 \geq 0.
\end{aligned}
\]

\textbf{(c)} Menetapkan variabel nonbasis $x = y = 0$ memberikan solusi basis
layak $(x, y, s_1, s_2) = (0, 0, 12, 24)$ dengan nilai fungsi objektif
$z = 0$.

\textbf{(d)} Berdasarkan pendakian tercuram, kita memilih koefisien objektif
terbesar: $5 > 4$, sehingga $y$ masuk terlebih dahulu.

\exsol{7.2}{Mengonversi ke Bentuk Standar}{1}

Fungsi objektif sudah berupa maksimisasi dan kedua variabel sudah
nonnegatif. Tambahkan variabel slack $s_1 \geq 0$ ke kendala $\leq$ dan
kurangkan variabel surplus $s_2 \geq 0$ dari kendala $\geq$:
\[
\begin{aligned}
\max \quad & 3x_1 - 5x_2 \\
\st \quad & 2x_1 + x_2 + s_1 = 6, \\
& -x_1 + 3x_2 - s_2 = 2, \\
& x_1, x_2, s_1, s_2 \geq 0.
\end{aligned}
\]

\exsol{7.3}{Mengonversi ke Bentuk Standar}{1}

Negasikan fungsi objektif untuk mengonversi min menjadi maks: $\max\; x_1 - 2x_2 - x_3$.
Karena $x_3$ bebas, tuliskan $x_3 = x_3^+ - x_3^-$ dengan
$x_3^+, x_3^- \geq 0$. Kendala pertama sudah berupa persamaan; tambahkan
slack $s_1 \geq 0$ ke kendala kedua:
\[
\begin{aligned}
\max \quad & x_1 - 2x_2 - x_3^+ + x_3^- \\
\st \quad & x_1 - x_2 + 2x_3^+ - 2x_3^- = 4, \\
& -2x_1 + x_2 + s_1 = 1, \\
& x_1, x_2, x_3^+, x_3^-, s_1 \geq 0.
\end{aligned}
\]

\exsol{7.4}{Mengonversi ke Bentuk Standar}{1}

Fungsi objektif sudah berupa maksimisasi. Tangani variabelnya: $x_1$
bebas, sehingga tetapkan $x_1 = x_1^+ - x_1^-$ dengan $x_1^+, x_1^- \geq 0$;
$x_2, x_3 \leq 0$, sehingga tetapkan $x_2 = -x_2'$ dan $x_3 = -x_3'$ dengan
$x_2', x_3' \geq 0$. Setelah substitusi, fungsi objektif menjadi
$5x_1^+ - 5x_1^- - x_2' + 3x_3'$. Kurangkan surplus $s_1 \geq 0$ dari
kendala $\geq$; persamaan tetap apa adanya:
\[
\begin{aligned}
\max \quad & 5x_1^+ - 5x_1^- - x_2' + 3x_3' \\
\st \quad & x_1^+ - x_1^- - 2x_2' + x_3' - s_1 = 7, \\
& -x_2' + 4x_3' = 3, \\
& x_1^+, x_1^-, x_2', x_3', s_1 \geq 0.
\end{aligned}
\]

\exsol{7.5}{Mengonversi ke Bentuk Standar}{1}

(Solusi Terpilih Buku.) Negasikan fungsi objektif:
$\max\; -2x_1 - 4x_2$. Kurangkan surplus $s_1 \geq 0$ dari kendala $\geq$
dan tambahkan slack $s_2 \geq 0$ ke kendala $\leq$. Karena
$x_1$ bebas, ganti $x_1 = x_1' - x_1''$ dengan
$x_1', x_1'' \geq 0$:
\[
\begin{aligned}
\max \quad & -2(x_1' - x_1'') - 4x_2 \\
\st \quad & (x_1' - x_1'') + x_2 - s_1 = 3, \\
& 2(x_1' - x_1'') - x_2 + s_2 = 5, \\
& x_1', x_1'', x_2, s_1, s_2 \geq 0.
\end{aligned}
\]

\exsol{7.6}{Mengonversi ke Bentuk Standar}{1}

(Solusi Terpilih Buku.) Negasikan fungsi objektif:
$\max\; -5y_1 + 3y_2 - y_3$. Tambahkan slack $s_1 \geq 0$ ke kendala $\leq$;
persamaan tidak memerlukan variabel baru. Karena $y_3$
bebas, ganti $y_3 = y_3' - y_3''$ dengan $y_3', y_3'' \geq 0$:
\[
\begin{aligned}
\max \quad & -5y_1 + 3y_2 - (y_3' - y_3'') \\
\st \quad & 2y_1 - y_2 + 3(y_3' - y_3'') = 7, \\
& -y_1 + 4y_2 + s_1 = 5, \\
& y_1, y_2, y_3', y_3'', s_1 \geq 0.
\end{aligned}
\]

\exsol{7.7}{Pivot ke arah yang berbeda}{2}

(Solusi Terpilih Buku, telah diverifikasi.) Memasukkan $x$ terlebih dahulu memerlukan satu
pivot tambahan dibandingkan lintasan pendakian tercuram: total tiga pivot.

\textbf{Pivot 1: $x$ masuk.} Uji rasio: $9/1 = 9$ ($s_1$),
$16/2 = 8$ ($s_2$), $14/1 = 14$ ($s_3$); minimumnya $8$, sehingga $s_2$
keluar. Menyelesaikan baris $s_2$ terhadap $x$ dan menyubstitusikannya:
\[
\begin{aligned}
\max\quad & z = 16 + 2y - s_2 \\
\st\quad & s_1 = 1 - \tfrac12 y + \tfrac12 s_2, \\
& x = 8 - \tfrac12 y - \tfrac12 s_2, \\
& s_3 = 6 - \tfrac32 y + \tfrac12 s_2.
\end{aligned}
\]

\textbf{Pivot 2: $y$ masuk} (satu-satunya koefisien positif). Rasio:
$1/\tfrac12 = 2$ ($s_1$), $8/\tfrac12 = 16$ ($x$), $6/\tfrac32 = 4$
($s_3$); sehingga $s_1$ keluar:
\[
\begin{aligned}
\max\quad & z = 20 - 4s_1 + s_2 \\
\st\quad & y = 2 - 2s_1 + s_2, \\
& x = 7 + s_1 - s_2, \\
& s_3 = 3 + 3s_1 - s_2.
\end{aligned}
\]

\textbf{Pivot 3: $s_2$ masuk} (koefisien $+1$). Baris $y$ memuat $s_2$
dengan koefisien positif, sehingga tidak memberikan batas; rasio $7/1 = 7$
($x$) dan $3/1 = 3$ ($s_3$), sehingga $s_3$ keluar:
\[
\begin{aligned}
\max\quad & z = 23 - s_1 - s_3 \\
\st\quad & x = 4 - 2s_1 + s_3, \\
& s_2 = 3 + 3s_1 - s_3, \\
& y = 5 + s_1 - s_3.
\end{aligned}
\]
Semua koefisien objektif negatif, sehingga kamus ini optimal:
$(x, y) = (4, 5)$ dengan $z^* = 23$, yaitu optimum yang sama yang dicapai lintasan
pendakian tercuram dalam dua pivot.

\exsol{7.8}{Latihan Metode Simpleks}{2}

(Solusi Terpilih Buku, telah diverifikasi.) Kamus awal:
\[
\begin{aligned}
\max\quad & z = 5x_1 + 3x_2 \\
\st\quad & s_1 = 12 - x_1 - x_2, \\
& s_2 = 16 - 2x_1 - x_2.
\end{aligned}
\]

\textbf{Pivot 1: $x_1$ masuk} (koefisien $5$). Rasio: $12/1 = 12$ dan
$16/2 = 8$, sehingga $s_2$ keluar:
\[
\begin{aligned}
\max\quad & z = 40 + \tfrac12 x_2 - \tfrac52 s_2 \\
\st\quad & x_1 = 8 - \tfrac12 x_2 - \tfrac12 s_2, \\
& s_1 = 4 - \tfrac12 x_2 + \tfrac12 s_2.
\end{aligned}
\]

\textbf{Pivot 2: $x_2$ masuk.} Rasio: $8/\tfrac12 = 16$ dan
$4/\tfrac12 = 8$, sehingga $s_1$ keluar:
\[
\begin{aligned}
\max\quad & z = 44 - s_1 - 2s_2 \\
\st\quad & x_2 = 8 - 2s_1 + s_2, \\
& x_1 = 4 + s_1 - s_2.
\end{aligned}
\]
Optimal: $(x_1, x_2) = (4, 8)$ dengan $z^* = 44$ (kedua kendala ketat).

\exsol{7.9}{Latihan Metode Simpleks}{2}

(Solusi Terpilih Buku, telah diverifikasi.) Kamus awal:
\[
\begin{aligned}
\max\quad & z = 5x_1 + 8x_2 \\
\st\quad & s_1 = 30 - x_1 - 2x_2, \\
& s_2 = 30 - 3x_1 - x_2.
\end{aligned}
\]

\textbf{Pivot 1: $x_2$ masuk} ($8 > 5$). Rasio: $30/2 = 15$ dan
$30/1 = 30$, sehingga $s_1$ keluar:
\[
\begin{aligned}
\max\quad & z = 120 + x_1 - 4s_1 \\
\st\quad & x_2 = 15 - \tfrac12 x_1 - \tfrac12 s_1, \\
& s_2 = 15 - \tfrac52 x_1 + \tfrac12 s_1.
\end{aligned}
\]

\textbf{Pivot 2: $x_1$ masuk.} Rasio: $15/\tfrac12 = 30$ dan
$15/\tfrac52 = 6$, sehingga $s_2$ keluar:
\[
\begin{aligned}
\max\quad & z = 126 - \tfrac{19}{5} s_1 - \tfrac25 s_2 \\
\st\quad & x_1 = 6 + \tfrac15 s_1 - \tfrac25 s_2, \\
& x_2 = 12 - \tfrac35 s_1 + \tfrac15 s_2.
\end{aligned}
\]
Optimal: $(x_1, x_2) = (6, 12)$ dengan $z^* = 126$.

\exsol{7.10}{Latihan Metode Simpleks}{2}

Tambahkan slack $s_1, s_2$:
\[
\begin{aligned}
\max\quad & z = 2x_1 + 3x_2 + x_3 \\
\st\quad & s_1 = 32 - 4x_1 - 2x_2 - 5x_3, \\
& s_2 = 28 - 2x_1 - 4x_2 - 3x_3.
\end{aligned}
\]

\textbf{Pivot 1: $x_2$ masuk} (koefisien terbesar $3$). Rasio:
$32/2 = 16$ dan $28/4 = 7$, sehingga $s_2$ keluar. Dengan
$x_2 = 7 - \tfrac12 x_1 - \tfrac34 x_3 - \tfrac14 s_2$:
\[
\begin{aligned}
\max\quad & z = 21 + \tfrac12 x_1 - \tfrac54 x_3 - \tfrac34 s_2 \\
\st\quad & s_1 = 18 - 3x_1 - \tfrac72 x_3 + \tfrac12 s_2, \\
& x_2 = 7 - \tfrac12 x_1 - \tfrac34 x_3 - \tfrac14 s_2.
\end{aligned}
\]

\textbf{Pivot 2: $x_1$ masuk} (koefisien $\tfrac12$). Rasio:
$18/3 = 6$ dan $7/\tfrac12 = 14$, sehingga $s_1$ keluar:
\[
\begin{aligned}
\max\quad & z = 24 - \tfrac{11}{6} x_3 - \tfrac16 s_1 - \tfrac23 s_2 \\
\st\quad & x_1 = 6 - \tfrac76 x_3 - \tfrac13 s_1 + \tfrac16 s_2, \\
& x_2 = 4 - \tfrac16 x_3 + \tfrac16 s_1 - \tfrac13 s_2.
\end{aligned}
\]
Semua koefisien objektif negatif, sehingga solusi ini optimal:
$(x_1, x_2, x_3) = (6, 4, 0)$ dengan $z^* = 24$; kedua kendala ketat
dan $x_3$ tidak pernah masuk.

\exsol{7.11}{Latihan Metode Simpleks}{2}

Tambahkan slack $s_1, s_2, s_3$:
\[
\begin{aligned}
\max\quad & z = 6x_1 + 8x_2 + 5x_3 \\
\st\quad & s_1 = 1800 - 4x_1 - x_2 - x_3, \\
& s_2 = 2000 - 2x_1 - 2x_2 - x_3, \\
& s_3 = 3200 - 4x_1 - 2x_2 - x_3.
\end{aligned}
\]

\textbf{Pivot 1: $x_2$ masuk} (koefisien $8$). Rasio:
$1800/1 = 1800$, $2000/2 = 1000$, $3200/2 = 1600$, sehingga $s_2$ keluar. Dengan
$x_2 = 1000 - x_1 - \tfrac12 x_3 - \tfrac12 s_2$:
\[
\begin{aligned}
\max\quad & z = 8000 - 2x_1 + x_3 - 4s_2 \\
\st\quad & s_1 = 800 - 3x_1 - \tfrac12 x_3 + \tfrac12 s_2, \\
& x_2 = 1000 - x_1 - \tfrac12 x_3 - \tfrac12 s_2, \\
& s_3 = 1200 - 2x_1 + s_2.
\end{aligned}
\]

\textbf{Pivot 2: $x_3$ masuk} (koefisien $+1$). Baris $s_3$ tidak
memuat $x_3$; rasio $800/\tfrac12 = 1600$ ($s_1$) dan
$1000/\tfrac12 = 2000$ ($x_2$), sehingga $s_1$ keluar:
\[
\begin{aligned}
\max\quad & z = 9600 - 8x_1 - 2s_1 - 3s_2 \\
\st\quad & x_3 = 1600 - 6x_1 - 2s_1 + s_2, \\
& x_2 = 200 + 2x_1 + s_1 - s_2, \\
& s_3 = 1200 - 2x_1 + s_2.
\end{aligned}
\]
Optimal: $(x_1, x_2, x_3) = (0, 200, 1600)$ dengan $z^* = 9600$; dua
kendala pertama ketat, kendala ketiga memiliki slack $s_3 = 1200$, dan $x_1$
tetap tidak masuk ke solusi optimal.

\exsol{7.12}{Latihan Metode Simpleks (Big-M)}{2}

Konversikan minimisasi menjadi maksimisasi, $\max\ z = -4x_1 - 6x_2 - 7x_3$,
kurangkan variabel surplus $e_1, e_2 \geq 0$, dan, karena titik asal
tidak layak, tambahkan variabel artifisial $a_1, a_2 \geq 0$ yang dikenai penalti $-M$
dalam fungsi objektif. Menyelesaikan terhadap $a_1, a_2$ memberikan kamus awal
\[
\begin{aligned}
\max\quad & z = -50M + (2M - 4)x_1 + (3M - 6)x_2 + (3M - 7)x_3 - M e_1 - M e_2 \\
\st\quad & a_1 = 20 - x_1 - x_2 - 2x_3 + e_1, \\
& a_2 = 30 - x_1 - 2x_2 - x_3 + e_2.
\end{aligned}
\]

\textbf{Pivot 1: $x_2$ masuk} (koefisien terbesar $3M - 6$). Rasio:
$20/1 = 20$ ($a_1$) dan $30/2 = 15$ ($a_2$), sehingga $a_2$ keluar:
\[
\begin{aligned}
\max\quad & z = (-90 - 5M) + \left(\tfrac{M}{2} - 1\right)x_1
  + \left(\tfrac{3M}{2} - 4\right)x_3 - M e_1
  + \left(\tfrac{M}{2} - 3\right)e_2 + \left(3 - \tfrac{3M}{2}\right)a_2 \\
\st\quad & a_1 = 5 - \tfrac12 x_1 - \tfrac32 x_3 + e_1 - \tfrac12 e_2 + \tfrac12 a_2, \\
& x_2 = 15 - \tfrac12 x_1 - \tfrac12 x_3 + \tfrac12 e_2 - \tfrac12 a_2.
\end{aligned}
\]

\textbf{Pivot 2: $x_3$ masuk} (koefisien $\tfrac{3M}{2} - 4$, yang
terbesar). Rasio: $5/\tfrac32 = \tfrac{10}{3}$ ($a_1$) dan
$15/\tfrac12 = 30$ ($x_2$), sehingga $a_1$ keluar. Kedua variabel artifisial sekarang
merupakan variabel nonbasis dan dapat dihapus:
\[
\begin{aligned}
\max\quad & z = -\tfrac{310}{3} + \tfrac13 x_1 - \tfrac83 e_1 - \tfrac53 e_2 \\
\st\quad & x_3 = \tfrac{10}{3} - \tfrac13 x_1 + \tfrac23 e_1 - \tfrac13 e_2, \\
& x_2 = \tfrac{40}{3} - \tfrac13 x_1 - \tfrac13 e_1 + \tfrac23 e_2.
\end{aligned}
\]

\textbf{Pivot 3: $x_1$ masuk} (koefisien $\tfrac13$). Rasio:
$\tfrac{10}{3}/\tfrac13 = 10$ ($x_3$) dan $\tfrac{40}{3}/\tfrac13 = 40$
($x_2$), sehingga $x_3$ keluar:
\[
\begin{aligned}
\max\quad & z = -100 - x_3 - 2e_1 - 2e_2 \\
\st\quad & x_1 = 10 - 3x_3 + 2e_1 - e_2, \\
& x_2 = 10 + x_3 - e_1 + e_2.
\end{aligned}
\]
Optimal: $(x_1, x_2, x_3) = (10, 10, 0)$ dengan $\max z = -100$, yaitu,
biaya minimum $z^* = 100$; kedua kendala $\geq$ ketat.

\begin{qaflag}
Catatan tentang kerumitan: dengan aturan masuk koefisien terbesar bawaan buku,
penyelesaian Big-M ini memerlukan tiga pivot dan melalui pecahan dengan penyebut tiga,
yang berada tepat di batas perhitungan manual. Terdapat rute yang lebih singkat:
memasukkan $x_1$ terlebih dahulu (koefisiennya $2M - 4$ juga positif) langsung mendorong
$a_1$ keluar dan mencapai optimum yang sama hanya dalam dua pivot.
Pengajar mungkin ingin menunjukkan lintasan tersebut kepada mahasiswa, atau buku dapat
menambahkan petunjuk seperti ``memasukkan $x_1$ terlebih dahulu mempersingkat perhitungan.''
Latihan itu sendiri benar sebagaimana dinyatakan.
\end{qaflag}

\exsol{7.13}{Latihan Metode Simpleks (Big-M)}{2}

Konversikan menjadi $\max\ z = -40x_1 - 48x_2 - 30x_3$, kurangkan surplus
$e_1, e_2 \geq 0$, dan tambahkan variabel artifisial berpenalti $a_1, a_2 \geq 0$:
\[
\begin{aligned}
\max\quad & z = -55M + (3M - 40)x_1 + (5M - 48)x_2 + (3M - 30)x_3 - M e_1 - M e_2 \\
\st\quad & a_1 = 25 - 2x_1 - 2x_2 - x_3 + e_1, \\
& a_2 = 30 - x_1 - 3x_2 - 2x_3 + e_2.
\end{aligned}
\]

\textbf{Pivot 1: $x_2$ masuk} (koefisien terbesar $5M - 48$). Rasio:
$25/2 = 12.5$ ($a_1$) dan $30/3 = 10$ ($a_2$), sehingga $a_2$ keluar:
\[
\begin{aligned}
\max\quad & z = (-480 - 5M) + \left(\tfrac{4M}{3} - 24\right)x_1
  + \left(2 - \tfrac{M}{3}\right)x_3 - M e_1
  + \left(\tfrac{2M}{3} - 16\right)e_2 + \left(16 - \tfrac{5M}{3}\right)a_2 \\
\st\quad & a_1 = 5 - \tfrac43 x_1 + \tfrac13 x_3 + e_1 - \tfrac23 e_2 + \tfrac23 a_2, \\
& x_2 = 10 - \tfrac13 x_1 - \tfrac23 x_3 + \tfrac13 e_2 - \tfrac13 a_2.
\end{aligned}
\]

\textbf{Pivot 2: $x_1$ masuk} ($\tfrac{4M}{3} - 24$ mengalahkan
$\tfrac{2M}{3} - 16$). Rasio: $5/\tfrac43 = \tfrac{15}{4}$ ($a_1$) dan
$10/\tfrac13 = 30$ ($x_2$), sehingga $a_1$ keluar. Kedua variabel artifisial sekarang
merupakan variabel nonbasis; setelah menghapusnya,
\[
\begin{aligned}
\max\quad & z = -570 - 4x_3 - 18e_1 - 4e_2 \\
\st\quad & x_1 = \tfrac{15}{4} + \tfrac14 x_3 + \tfrac34 e_1 - \tfrac12 e_2, \\
& x_2 = \tfrac{35}{4} - \tfrac34 x_3 - \tfrac14 e_1 + \tfrac12 e_2.
\end{aligned}
\]
Semua koefisien negatif, sehingga solusi ini optimal hanya setelah dua pivot:
$(x_1, x_2, x_3) = \left(\tfrac{15}{4}, \tfrac{35}{4}, 0\right)$ dengan
biaya minimum $z^* = 40 \cdot \tfrac{15}{4} + 48 \cdot \tfrac{35}{4}
= 150 + 420 = 570$. Kedua kendala $\geq$ ketat.

\begin{qaflag}
Keputusan yang perlu dipertimbangkan: solusi optimal berbentuk pecahan
$\left(\tfrac{15}{4}, \tfrac{35}{4}\right)$ meskipun semua datanya
bilangan bulat. Penyelesaiannya singkat (dua pivot) dan nilai optimal $570$
rapi, sehingga ini dapat diterima, tetapi jika jawaban bilangan bulat lebih disukai,
mengubah ruas kanan pertama dari $25$ menjadi $24$ menghasilkan
optimum bulat $(3, 9, 0)$ dengan biaya $552$ (telah diverifikasi).
\end{qaflag}

\exsol{7.14}{Optimalitas Kamus}{2}

(Solusi Terpilih Buku, telah diverifikasi.) Solusi basis dibaca dari
suku konstanta: $x_2 = 1$, $x_4 = 2$, $x_3 = b$, dengan
$x_1 = x_5 = x_6 = 0$.
\begin{enumerate}
\item \textbf{Layak dan optimal:} kelayakan mengharuskan setiap suku konstanta
nonnegatif, sehingga $b \geq 0$; optimalitas mengharuskan tidak ada arah
perbaikan, sehingga $c_1 \leq 0$, $c_5 \leq 0$, $c_6 \leq 0$. Nilai
$a_1, a_2$ boleh sembarang.
\item \textbf{Tidak layak:} $b < 0$ (konstanta $1$ dan $2$
sudah nonnegatif, sehingga hanya baris $x_3$ yang dapat gagal).
\item \textbf{Degenerat:} $b = 0$, sehingga variabel basis $x_3$
bernilai nol (kelayakan bagian lainnya otomatis).
\item \textbf{Layak dan tak terbatas:} kita memerlukan $b \geq 0$ bersama suatu
variabel masuk yang koefisien objektifnya positif dan peningkatannya
tidak pernah menurunkan variabel basis apa pun. Variabel $x_5$ muncul
dengan koefisien tetap $-1$ pada baris $x_3$ dan $x_6$ dengan $-1$ pada kedua
baris $x_4$ dan $x_3$, sehingga uji rasio selalu menghentikannya. Hanya $x_1$
yang dapat membuktikan ketakterbatasan: koefisien kolomnya adalah $+1$, $+1$, dan
$-a_2$, sehingga kondisinya adalah $c_1 > 0$ dan $a_2 \leq 0$ (dengan $b \geq 0$
untuk kelayakan).
\end{enumerate}

\exsol{7.15}{Solusi Basis Layak}{2}

\begin{qaflag}
Perbaikan buku telah diterapkan: dalam sumber, semua kecuali satu kendala latihan ini
telah dikomentari, dan satu kendala yang terlihat
($-x_1 + x_2 \leq 2$) tidak sesuai dengan gambar. Sistem kendala
direkonstruksi tepat dari gambar (kelima titik sudut berlabel semuanya cocok:
$A = (2,1)$, $B = (5,1)$, $C = (7,3)$, $D = (0, \tfrac{16}{3})$,
$E = (0, \tfrac{7}{3})$) dan pernyataan buku diperbaiki menjadi
\[
\max\ 2x_1 + x_2 \quad \st\quad 2x_1 + 3x_2 \geq 7,\ \
x_2 \geq 1,\ \ x_1 - x_2 \leq 4,\ \ x_1 + 3x_2 \leq 16,\ \ x_1, x_2 \geq 0,
\]
dengan variabel surplus $s_1, s_2$ pada kedua kendala $\geq$ dan slack
$s_3, s_4$ pada kedua kendala $\leq$. Solusi di bawah ini sesuai dengan
pernyataan yang telah diperbaiki.
\end{qaflag}

\textbf{(a)} Titik sudut $D = (0, \tfrac{16}{3})$ adalah perpotongan
batas $x_1 = 0$ dan $x_1 + 3x_2 = 16$ (yaitu, $s_4 = 0$). Jadi,
variabel nonbasis adalah $N = \{x_1, s_4\}$ dan variabel basis adalah
$B = \{x_2, s_1, s_2, s_3\}$, dengan nilai
\[
x_2 = \tfrac{16}{3}, \qquad s_1 = 9, \qquad s_2 = \tfrac{13}{3},
\qquad s_3 = \tfrac{28}{3}.
\]

\textbf{(b)} Bergerak dari $D$ ke $C = (7, 3)$ berarti menempuh sisi
$s_4 = 0$ ketika $x_1$ meningkat dari $0$: jadi \emph{$x_1$ masuk}. Pada $C$,
batas $x_1 - x_2 = 4$ menjadi ketat, yaitu, $s_3 = 0$: jadi
\emph{$s_3$ keluar}. Ini sesuai dengan uji rasio dalam kamus di
$D$, ketika peningkatan $x_1$ (dengan $s_4 = 0$) mendorong
$s_3 = \tfrac{28}{3} - \tfrac43 x_1$ ke nol terlebih dahulu, pada $x_1 = 7$,
sebelum $s_2$ (pada $x_1 = 13$) atau $x_2$ (pada $x_1 = 16$).

\exsol{7.16}{Apa yang dilindungi oleh uji rasio}{2}

\textbf{(a)} Melakukan pivot pada baris $s_2$ berarti menyelesaikan
$s_2 = 9 - \tfrac32 x + \tfrac12 s_3$ terhadap $x$, sehingga memberikan
$x = 6 - \tfrac23 s_2 + \tfrac13 s_3$, lalu menyubstitusikannya ke semua tempat:
\[
\begin{aligned}
\max\quad & z = 24 - \tfrac13 s_2 - \tfrac43 s_3 \\
\st\quad & s_1 = -1 + \tfrac13 s_2 + \tfrac13 s_3, \\
& x = 6 - \tfrac23 s_2 + \tfrac13 s_3, \\
& y = 4 + \tfrac13 s_2 - \tfrac23 s_3.
\end{aligned}
\]
Solusi basisnya adalah $(x, y, s_1, s_2, s_3) = (6, 4, -1, 0, 0)$: variabel
basis $s_1$ negatif, sehingga kamus tidak layak. (Memang,
titik $(x, y) = (6, 4)$ melanggar $x + y \leq 9$.) Perhatikan bahwa
``nilai fungsi objektif'' $24$ melebihi optimum sebenarnya $23$ justru karena
titik tersebut berada di luar daerah layak.

\textbf{(b)} Ketika variabel masuk $x$ meningkat, setiap variabel basis
dengan koefisien $x$ negatif menurun menuju nol; rasio setiap
baris mengukur seberapa jauh $x$ dapat bergerak sebelum variabel basis pada baris itu mencapai
nol. Rasio minimum adalah nilai pemblokir pertama. Melakukan pivot pada
baris lain mendorong $x$ melampaui pemblokir pertama tersebut, sehingga variabel basis yang
memiliki rasio minimum (di sini $s_1$, terblokir pada $x = 4$) terdorong
menjadi negatif. Aturan rasio minimum tepat melindungi nonnegativitas
semua variabel basis yang tetap berada dalam basis.

\exsol{7.17}{Hasil imbang dalam uji rasio memaksa degenerasi}{3}

(Solusi Terpilih Buku, telah diverifikasi, dengan bagian (c) dinyatakan ulang.)

\textbf{(a)} Kamus awalnya adalah
\[
\begin{aligned}
\max\quad & z = 2x + y \\
\st\quad & s_1 = 4 - x - y, \\
& s_2 = 8 - 2x - y.
\end{aligned}
\]
Dengan $x$ masuk (dan $y = 0$), uji rasio memberikan $4/1 = 4$ untuk $s_1$
dan $8/2 = 4$ untuk $s_2$: terjadi hasil imbang pada $x = 4$.

\textbf{(b)} Biarkan $s_1$ keluar. Maka $x = 4 - y - s_1$ dan
\[
\begin{aligned}
\max\quad & z = 8 - y - 2s_1 \\
\st\quad & x = 4 - y - s_1, \\
& s_2 = 0 + y + 2s_1.
\end{aligned}
\]
Variabel basis $s_2$ bernilai $0$: kamus tersebut degenerat. Kamus ini
juga optimal (kedua biaya tereduksi negatif), dengan $(x, y) = (4, 0)$ dan
$z^* = 8$.

\textbf{(c)} Sebagai gantinya, biarkan $s_2$ keluar. Maka
$x = 4 - \tfrac12 y - \tfrac12 s_2$ dan
\[
\begin{aligned}
\max\quad & z = 8 + 0 \cdot y - s_2 \\
\st\quad & x = 4 - \tfrac12 y - \tfrac12 s_2, \\
& s_1 = 0 - \tfrac12 y + \tfrac12 s_2.
\end{aligned}
\]
Sekali lagi, suatu variabel basis ($s_1$) sama dengan nol: degenerat. Biaya tereduksi
$y$ tepat $0$. Biaya tereduksi nol biasanya menandakan optimum alternatif
sepanjang sisi yang bersesuaian; di sini baris degenerat
$s_1 = 0 - \tfrac12 y + \cdots$ langsung memblokir $y$, sehingga
\emph{solusi} optimal $(4, 0)$ unik, tetapi terdapat beberapa
\emph{basis} optimal yang mendeskripsikannya. Ini merupakan ciri khas degenerasi.

\textbf{(d)} Jika dua baris memiliki rasio minimum yang sama, variabel masuk
berhenti pada nilai ketika kedua variabel basis yang bersesuaian mencapai nol
secara bersamaan. Hanya salah satunya yang keluar dari basis; yang lain tetap
menjadi variabel basis dengan nilai nol, sehingga kamus baru pasti memiliki variabel basis
yang sama dengan nol, yang merupakan definisi degenerasi. Sebuah
konstruksi: $\max\ x + y$ dengan kendala $x \leq 3$, $x + 2y \leq 3$,
$x, y \geq 0$; memasukkan $x$ memberikan rasio $3/1 = 3$ dan $3/1 = 3$, yaitu hasil imbang.
